\subsection{The Mayor Agent: Macro-Level Intervention and DRL Execution}

The \texttt{MayorAgent} serves as the centralized executive entity within the simulation. Unlike the \texttt{BarangayAgents}, which operate continuously on static, hardcoded baseline budgets, the Mayor Agent is fully dynamic and serves as the physical actuator for the Heuristic-Guided Deep Reinforcement Learning (HuDRL) policy. To systematically detail its computational role, this section is divided into four parts: (1) The temporal decision epoch, (2) The budget actuation pipelines, (3) The political survival constraint, and (4) The closed-loop communication topology.

\paragraph{Part 1: The Macro-Level Decision Epoch (The AI Bridge)}

Computationally, the \texttt{MayorAgent} acts as the structural bridge between the simulated Agent-Based Model (ABM) environment and the external PyTorch-based neural network. To simulate the administrative realities of municipal governance, the agent does not evaluate policy daily. Instead, its \texttt{step()} function enforces a strict 90-day (quarterly) temporal epoch \citep{Akbarpour2021}. 

At the start of each quarter ($t \pmod{90} == 0$), the Mayor Agent pauses local micro-executions to observe the global state vector (aggregated compliance, average attitudes, remaining annual budget, and political capital). It queries the HuDRL algorithm, receiving a continuous action vector representing a normalized budget allocation strategy. Once received, the Mayor Agent immediately translates these abstract floating-point values into physical municipal interventions.

\paragraph{Part 2: Budget Translation and Macro-Intervention Levers}
The AI's action vector is partitioned and executed across the same three policy levers utilized by the Barangays, but at a massively scaled municipal capacity \citep{Ibanez2022}:

\begin{enumerate}
    \item \textbf{Macro-IEC (Mass Media Campaigns):} When the AI allocates funds to a barangay's IEC sector, the Mayor bypasses grassroots methods and invests in mass media. The allocation ($B_{lgu\_iec}$) funds broad-reach radio broadcasts ($C_{rad} = \text{\textpeso} 500$) and large-scale assembly events ($C_{evt} = \text{\textpeso} 5,000$). The intensity of municipal IEC ($I_{macro}$) is bounded by target saturation levels:
    \begin{align}
        n_{rad} &= \min \left( T_{rad}, \left\lfloor \frac{B_{lgu\_iec}}{C_{rad}} \right\rfloor \right) \\
        rem_{lgu} &= \max \left( 0, B_{lgu\_iec} - (n_{rad} \times C_{rad}) \right) \\
        n_{evt} &= \min \left( T_{evt}, \left\lfloor \frac{rem_{lgu}}{C_{evt}} \right\rfloor \right) \\
        I_{macro} &= \min \left( 1.0, \omega_{rad} \left( \frac{n_{rad}}{T_{rad}} \right) + \omega_{evt} \left( \frac{n_{evt}}{T_{evt}} \right) \right)
        \label{eq:macro_iec_logic}
    \end{align}
    \addequation{Municipal IEC Intensity and Resource Partitioning}

    \noindent \textit{Where:}
    \begin{itemize}
        \item $B_{lgu\_iec}$: The municipal budget allocation injected by the AI Mayor.
        \item $n_{rad}$ and $n_{evt}$: The number of radio broadcasts and assembly events funded.
        \item $rem_{lgu}$: The residual municipal budget used in the allocation cascade.
        \item $I_{macro}$: The resulting Macro-Level IEC Intensity scalar.
        \item $T_{rad}$ and $T_{evt}$: Target saturation limits for broadcasts and events.
    \end{itemize}

    \item \textbf{Macro-Enforcement (LGU Strike Force) and Object Destruction:} To execute localized crackdowns, enforcement funds are translated into elite ``Strike Force'' personnel. Deploying a single municipal enforcer on a strict 30-day temporal contract consumes exactly \textpeso 12,000 (based on Region X minimum wage). The municipal enforcement cost ($C_{\mathrm{Enf}}$) is linearly tied to human capital:
    \begin{equation}
        C_{\mathrm{Enf}} = (N_{\mathrm{Enforcers}} \times W_{\mathrm{Daily}} \times 30)
        \label{eq:macro_enf_cost}
    \end{equation}
    \addequation{Municipal Strike Force Deployment Cost Calculation}
    
    The Mayor dynamically instantiates these high-capacity \texttt{EnforcementAgents} (flagged as \texttt{is\_municipal = True}). They achieve a 100\% daily patrol rate, execute up to 25 apprehensions per tick, and issue severe \textpeso 1,000 citations \citep{Wang2023}. Crucially, these agents are bound by their strict 30-day limit enforced via computational object destruction. Upon reaching zero, the object invokes a self-destruct sequence (\texttt{self.model.grid.remove\_agent()}), physically deleting itself from the spatial grid. 
    
    \item \textbf{Macro-Incentives (Economic Augmentation):} Funds allocated to positive reinforcement are channeled directly into the target \texttt{BarangayAgent}'s Tier 2 financial ledger, ensuring liquidity for citizen reward claims.
\end{enumerate}

\paragraph{Part 3: Political Capital Maintenance (Survival Mechanics)}

To prevent the AI from adopting a computationally optimal but socially disastrous draconian strategy, the Mayor Agent is bound by a ``Political Capital'' constraint \citep{Hertweck2023}. Every municipal fine (exceeding \textpeso 500) levied by the Strike Force triggers an interrupt that applies a fractional deduction to the Mayor's public approval rating:
\begin{equation}
    \Delta PC = -0.00010 \times N_{fines}
    \label{eq:political_capital_decay}
\end{equation}
\addequation{Political Capital Decay via Punitive Enforcement}

If the Mayor's Political Capital drops below the critical threshold of $0.10$, the \texttt{BacolodModel} forces a premature termination of the simulation due to ``Political Collapse.''

\paragraph{Part 4: Inter-Agent Communication Topology}

The resulting multi-level architecture forms a closed, bi-directional computational feedback loop, officially grounding the ABM as a Markov Decision Process (MDP) \citep{Kompella2020}:
\begin{itemize}
    \item \textbf{Bottom-Up Aggregation:} Barangay Agents aggregate individual states into localized metrics ($R_{comp}$) and pass them to the centralized model.
    \item \textbf{Top-Down Actuation:} The Mayor Agent consults the external DRL policy and deploys financial capital downward to alter the grid or mutate household variables.
\end{itemize}

\vspace{0.5cm} 
\begin{algorithm}[H]
\caption{The Mayor Agent's Quarterly Intervention Cycle}
\label{alg:mayor_decision_english}
\begin{algorithmic}[1]
\State \textbf{Starting Information Needed:} 
\State \parbox[t]{0.9\linewidth}{$\rightarrow$ The current global segregation rate, remaining annual budget, and political approval rating.\strut}
\Statex
\State \textbf{Step 1: The AI Consultation (Occurs every 90 Days)}
\State \parbox[t]{0.9\linewidth}{$\rightarrow$ Pause the simulation and transmit the city's current status to the HuDRL Artificial Intelligence.\strut}
\State \parbox[t]{0.9\linewidth}{$\rightarrow$ Receive a finalized, mathematically optimized budget blueprint.\strut}
\Statex
\State \textbf{Step 2: Execute Budget Actuations (Top-Down Flow)}
\For{every Barangay in the municipality}
    \State \parbox[t]{0.88\linewidth}{$\rightarrow$ \textit{Information Campaigns:} Purchase radio airtime and fund large events to boost household motivation.\strut}
    \State \parbox[t]{0.88\linewidth}{$\rightarrow$ \textit{Strike Force Deployment:} Spawn elite inspectors into the grid on strict 30-day contracts.\strut}
    \State \parbox[t]{0.88\linewidth}{$\rightarrow$ \textit{Incentive Augmentation:} Subsidize local Barangay ledgers for citizen payouts.\strut}
\EndFor
\Statex
\State \textbf{Step 3: Monitor Survival Constraints}
\If{the Mayor's Political Capital $\le 10\%$}
    \State \parbox[t]{0.88\linewidth}{$\rightarrow$ Trigger an immediate simulation halt due to administrative collapse.\strut}
\Else
    \State \parbox[t]{0.88\linewidth}{$\rightarrow$ Resume the simulation for the next 90-day decision epoch.\strut}
\EndIf
\end{algorithmic}
\end{algorithm}